\section{Parametric specification}

\begin{figure*}[ht!]
  %\renewcommand\thesubfigure{\arabic{subfigure}}
  \centering
  \begin{subfigure}[t]{0.11\textwidth}
    \includegraphics[width=\textwidth]{1-evaluation-square}
    \caption{}
    \label{fig:1-eval}
  \end{subfigure}
  \hfill
  \begin{subfigure}[t]{0.11\textwidth}
    \includegraphics[width=\textwidth]{2-evaluation-extrusion}
    \caption{}
    \label{fig:2-eval}
  \end{subfigure}
  \hfill
  \begin{subfigure}[t]{0.11\textwidth}
    \includegraphics[width=\textwidth]{3-evaluation-insert}
    \caption{}
    \label{fig:3-eval}
  \end{subfigure}
  \hfill
  \begin{subfigure}[t]{0.11\textwidth}
    \includegraphics[width=\textwidth]{4-evaluation-extrusion}
    \caption{}
    \label{fig:4-eval}
  \end{subfigure}
  \hfill
  \begin{subfigure}[t]{0.11\textwidth}
    \includegraphics[width=\textwidth]{5-evaluation-triangulation}
    \caption{}
    \label{fig:5-eval}
  \end{subfigure}
  \hfill
  \begin{subfigure}[t]{0.11\textwidth}
    \includegraphics[width=\textwidth]{6-evaluation-collapse}
    \caption{}
    \label{fig:6-eval}
  \end{subfigure}
  \hfill
  \begin{subfigure}[t]{0.11\textwidth}
    \includegraphics[width=\textwidth]{7-evaluation-chamfer}
    \caption{}
    \label{fig:7-eval}
  \end{subfigure}
  \hfill
  \begin{subfigure}[t]{0.11\textwidth}
    \includegraphics[width=\textwidth]{8-evaluation-colour}
    \caption{}
    \label{fig:8-eval}
  \end{subfigure}
  \caption{Evaluation. %'s applications. 
  (a) $1$-square(pos); (b)
    $2$-extrude($PN_{1}$,vec); (c) $3$-insert($PN_{2}$); (d)
    $4$-extrude($PN_{3}$,vec); (e) $5$-triangulate($PN_{4}$); (f)
    $6$-collapse($PN_{5}$); (g) $7$-chamfer($PN_{6}$); (h) $8$-colour($PN_{7}$)}
  \label{fig:evaluation}
\end{figure*}

\begin{figure*}[ht!]
  %\renewcommand\thesubfigure{\arabic{subfigure}}
  \centering
  \begin{subfigure}[t]{0.11\textwidth}
    \includegraphics[width=\textwidth]{1-reevaluation-square}
    \caption{}
    \label{fig:1-reeval}
  \end{subfigure}
  \hfill
  \begin{subfigure}[t]{0.11\textwidth}
    \includegraphics[width=\textwidth]{1ADD-reevaluation-insert}
    \caption{}
    \label{fig:1add-reeval}
  \end{subfigure}
  \hfill
  \begin{subfigure}[t]{0.11\textwidth}
    \includegraphics[width=\textwidth]{2-reevaluation-extrusion}
    \caption{}
    \label{fig:2-reeval}
  \end{subfigure}
  \begin{subfigure}[t]{0\textwidth}
    
\includegraphics[width=\textwidth]{3-reevaluation-insert}
      \caption{}
    \label{fig:3delete-reeval}
  \end{subfigure}
  \hfill
  \begin{subfigure}[t]{0.11\textwidth}
    \includegraphics[width=\textwidth]{4-reevaluation-extrusion}
    \caption{}
    \label{fig:4-reeval}
  \end{subfigure}
  \hfill
  \begin{subfigure}[t]{0.11\textwidth}
    \includegraphics[width=\textwidth]{5-reevaluation-triangulation}
    \caption{}
    \label{fig:5-reeval}
  \end{subfigure}
  \hfill
  \begin{subfigure}[t]{0.11\textwidth}
    \includegraphics[width=\textwidth]{6-reevaluation-collapse}
    \caption{}
    \label{fig:6-reeval}
  \end{subfigure}
  \hfill
  \begin{subfigure}[t]{0.11\textwidth}
    \includegraphics[width=\textwidth]{7-reevaluation-chamfer}
    \caption{}
    \label{fig:7-reeval}
  \end{subfigure}
  \hfill
  \begin{subfigure}[t]{0.11\textwidth}
    \includegraphics[width=\textwidth]{8-reevaluation-colour}
    \caption{}
    \label{fig:8-reeval}
  \end{subfigure}
  \caption{Re-evaluation. (a) $1$-square(pos); (b) ADD$1$-insert(3); (c)
    $2$-extrude($PN_{1}$,vec); (d) DELETE $3$-insert($PN_{2}$); (e)
    $4$-extrude($PN_{3}$,vec); (f) $5$-triangulate($PN_{4}$); (g)
    $6$-collapse($PN_{5}$); (h) $7$-chamfer($PN_{6}$); (i) $8$-colour($PN_{7}$)}
  \label{fig:reevaluation}
\end{figure*}

During an object's construction, a \emph{parametric specification} records both the rules
representing the applied operations
%applied 
and their pa\-ra\-me\-ters (both topological and geometric) in order to
describe the modeling process. %
%When a parametric specification is edited, 
Editing a parametric specification means that
rules may be added, deleted, moved,
and their parameters can be changed. %
With such changes, the topological parameters may have different concrete names,
be deleted and so on. %
As a result, using the parameters' concrete names eventually lead to unexpected
results (at best) or failure at runtime. %
To this end, persistent names are required to robustly identify topological
parameters across an object re-evaluation. %

This section will follow the evaluation (Fig.~\ref{fig:evaluation}) of a modeling
process as an example in order to illustrate the creation of a parametric
specification and its persistent names. %
Its re-evaluation (Fig.~\ref{fig:reevaluation}) illustrates a variant modeling
process where a rule is added between applications 1 and 2, and where the
application 3 is deleted. %

%All the figures in this section and the next ones were generated using Jerboa and the development presented in this paper. %

In this paper, all the figures were generated using Jerboa and the software overlay developed to implement the concepts presented in this section and the following ones. An overview of the results can be found here: \url{http://xlim-sic.labo.univ-poitiers.fr/jerboa/doc/model-reevaluation-based-on-graph-transformation-rules/}

\subsection{Persistent name}

Since rules use darts as topological parameters, it follows that each persistent
name must represent a unique dart. %
As it happens, rules make it possible to determine unambiguously, which node filters or creates any dart. %

\begin{figure}
    \centering
    \includegraphics[height=.25\textwidth]{figures/rule-create-square-face.pdf}
    \caption{Square creation rule}
    \label{fig:rule-create-square}
\end{figure}

For example, the square rule (Fig.~\ref{fig:rule-create-square}) creates $8$
darts \emph{ex~nihilo} (Fig.~\ref{fig:1-eval}), one per node. %
Thus, dart $3$ is created during the first application, by the square's node $n_{3}$. The history of dart $3$ is
noted $[1n_{3}]$. %
Similarly, since dart $5$ (resp.) $6$ is created by node $n_{5}$ (resp. $n_{6}$), its history is $[1n_{5}]$ (resp. $[1n_{6}]$). %

The second application (Fig.~\ref{fig:2-eval} and for more details
Fig.~\ref{fig:before-after-extrusion}) of the
face extrusion rule (Fig.~\ref{fig:rule-face-extrusion}), creates darts from the
matched darts. %
This rule is applied on the square face created during the first application. Since this face contains 8 darts, each of them is matched by node $n1$ of the extrusion rule (Fig~\ref{fig:before-after-extrusion}). %
As a con\-se\-quen\-ce, nodes $n_{2}$ to $n_{6}$ create darts copied from the
8 matched darts. %
Therefore, dart $6$'s history is now $[1n_{6};2n_{1}]$, its copy $8$'s
history is $[1n_{6};2n_{2}]$, its copy $9$'s history is $[1n_{6};2n_{3}]$, its
copy $10$'s history is $[1n_{6};2n_{4}]$ and so on. %
Similarly, histories of darts $46$ and $47$ are $[1n_{5};2n_{5}]$ and $[1n_{5};2n_{6}]$, respectively.

For the same reason, histories of darts $33$ and $35$  are $[1n_{3};2n_{2}]$ and
$[1n_{3};2n_{4}]$, respectively. %
It follows that the cube's 48 darts all have a different history. %

\begin{figure}[ht]
  \centering
  \includegraphics[height=30pt]{rule-insert-vertex}
  \caption{Rule defining the insertion of a vertex on an edge}
  \label{fig:rule-insert-vertex}
\end{figure}

The third application (Fig.~\ref{fig:3-eval}) uses the rule inserting a vertex on an edge.
Both darts $46$ and $47$ are matched by  node $n_{0}$ of the insertion rule. %
Their respective histories are $[1n_{5};2n_{5};3n_{0}]$ and $[1n_{5};2n_{6};3n_{0}]$ now.
Conversely, darts $33$ and $35$ have not been not matched during the vertex insertion and
their respective histories remain $[1n_{3};2n_{2}]$ and $[1n_{3};2n_{4}]$. %
The object's 52 darts, again, have each a different history. %

In short, the history of any dart of the topological model is en\-ti\-re\-ly defined by the rules applied during the evaluation process.
% The same goes for the next steps’ darts in this or any other construction made by successive applications of rules
This process guarantees to associate each dart with a unique his\-to\-ry; hence, we use this history as the persistent name of the dart.
Therefore, $PN_{1}=[1n_{6}]$, $PN_{2}=[1n_{5};2n_{5}]$,
$PN_{3}=[1n_{5};2n_{6};3n_{0}]$ and so on. %

\subsection{Parametric Specification syntax}

A number of fields are used to describe an application within a parametric
specification, namely an application number, a rule ap\-pli\-ca\-tion with 
the topological pa\-ra\-me\-ters and the geometric parameters. %
\begin{figure}
    \centering
\small
\begin{verbatim}
1-square(pos)
2-extrude(PN1=[1n6], vec)
3-insert(PN2=[1n5;2n5])
4-extrude(PN3=[1n5;2n6;3n0], vec)
5-triangulate(PN4=[1n3;2n2])
6-collapse(PN5=[1n4;2n6;4n6])
7-chamfer(PN6=[1n5;2n6;3n0;4n4;6n2])
8-colour(PN7=[1n3;2n4;5n0])
\end{verbatim}
    \caption{Initial parametric specification}
    \label{fig:init-spec}
\end{figure}
The parametric specification above represents the modeling process of an
initial evaluation (Fig.~\ref{fig:evaluation}) where each persistent name 
uses the history of the dart's number displayed in the previous construction
step (dart $6$ for $PN_1$, $46$ for $PN_2$, $47$ for $PN_3$, $33$ for $PN_4$
and so on.). %

Furthermore, a set of tags are used to describe an application whenever it is either
added (\texttt{ADD}), deleted (\texttt{DELETE}) or moved (\texttt{MOVE}) at
reevaluation. %
Thus, the initial parametric specification, once edited as %above, 
shown in Fig.~\ref{fig:reval-spec}
produces the
reevaluation process 
shown in Fig.~\ref{fig:reevaluation}.
\begin{figure}
    \centering
    \small
\begin{verbatim}
1-square(pos)
ADD1-insert(3)
2-extrude(PN1=[1n6], vec)
DELETE 3-insert(PN2=[1n5;2n4])
...
8-colour(PN7=[1n3;2n3;5n0])
\end{verbatim}
    \caption{Edited parametric specification}
    \label{fig:reval-spec}
\end{figure}

%%% Local Variables:
%%% mode: latex
%%% TeX-master: "../main"
%%% End:
